\section{Methodology}
\label{sec:methodology}

We present the Faithfulness-Hallucination Index (FHI), a composite metric that quantifies the relationship between explanation faithfulness and hallucination propensity. The framework operates in five stages: (1) answer and explanation generation, (2) attribution analysis, (3) causal perturbation, (4) multi-dimensional metric computation, and (5) composite scoring. Figure~\ref{fig:pipeline} illustrates the overall pipeline.

\subsection{Problem Formulation}

Given a question $q$, an LLM $\mathcal{M}$ generates an answer $a$ and an explanation $e$:
\begin{equation}
    (a, e) = \mathcal{M}(q)
\end{equation}

We further extract internal representations: token log-probabilities $\mathbf{p} = \{p_1, \ldots, p_n\}$ and attention matrices $\mathbf{A} = \{A^{(1)}, \ldots, A^{(L)}\}$ across $L$ layers. The goal is to compute $\text{FHI}(q, a, e, \mathbf{p}, \mathbf{A}) \in [0, 1]$, where high values indicate faithful reasoning and low values indicate hallucination risk.

\subsection{Attribution Analysis}

We employ attention rollout~\cite{abnar2020quantifying} as the primary attribution method. For a transformer with $L$ layers, the rollout matrix $\hat{A}$ is computed recursively:
\begin{equation}
    \hat{A}^{(l)} = \begin{cases}
        A^{(1)} & \text{if } l = 1 \\
        \left(\frac{A^{(l)} + I}{2}\right) \cdot \hat{A}^{(l-1)} & \text{otherwise}
    \end{cases}
\end{equation}
where $I$ is the identity matrix accounting for residual connections. The final attribution vector $\boldsymbol{\alpha} = \hat{A}^{(L)}[{-1}, :]$ represents the aggregated attention from the last generated token to all input tokens. We select the top-$K$ tokens by attribution score as the \textit{attribution set} $\mathcal{T}_\alpha$.

\subsection{Metric 1: Attribution Alignment Score (AAS)}

AAS measures the semantic overlap between the explanation tokens and the attribution-highlighted tokens. Let $\mathcal{T}_e$ denote the set of content tokens in the explanation $e$ (after stopword removal), and $\mathcal{T}_\alpha$ the top-$K$ attributed tokens. We compute:
\begin{equation}
    \text{AAS} = \frac{|\mathcal{T}_e \cap \mathcal{T}_\alpha|}{|\mathcal{T}_e \cup \mathcal{T}_\alpha|}
\end{equation}

This Jaccard similarity measures whether the model's explanation references the same tokens that its attention mechanism actually weighted. When multiple attribution methods are available, we average across methods:
\begin{equation}
    \text{AAS}_{\text{avg}} = \frac{1}{M} \sum_{m=1}^{M} \text{AAS}_m
\end{equation}

\subsection{Metric 2: Causal Impact Score (CIS)}

CIS quantifies whether explanation-highlighted tokens causally influence the model's output. For each perturbation strategy $s \in \{\text{mask}, \text{delete}, \text{replace}\}$, we:

\begin{enumerate}
    \item Construct a perturbed prompt $q'_s$ by applying strategy $s$ to the top-$K$ attributed tokens in the original question.
    \item Generate a new answer: $a'_s = \mathcal{M}(q'_s)$.
    \item Compute the output shift $\Delta_s$ as:
\end{enumerate}
\begin{equation}
    \Delta_s = \frac{1}{3}\left(D_{\text{JS}}(\mathbf{p}, \mathbf{p}'_s) + d_{\text{sem}}(a, a'_s) + |\text{ROUGE-L}(a, g) - \text{ROUGE-L}(a'_s, g)|\right)
\end{equation}
where $D_{\text{JS}}$ is the Jensen-Shannon divergence between log-probability distributions, $d_{\text{sem}}$ is the cosine distance between sentence embeddings (using all-MiniLM-L6-v2~\cite{reimers2019sentence}), and $g$ is the gold answer.

The final CIS is the mean shift across strategies:
\begin{equation}
    \text{CIS} = \frac{1}{|S|} \sum_{s \in S} \Delta_s
\end{equation}

High CIS indicates that the attributed tokens are causally important---removing them significantly changes the output.

\subsection{Metric 3: Explanation Stability Score (ESS)}

ESS measures the consistency of explanations across $N$ independent generations using the same prompt:
\begin{equation}
    \text{ESS} = \frac{2}{N(N-1)} \sum_{i=1}^{N} \sum_{j=i+1}^{N} \text{sim}_{\text{sem}}(e_i, e_j)
\end{equation}
where $\text{sim}_{\text{sem}}$ denotes cosine similarity between sentence embeddings of explanations $e_i$ and $e_j$. ESS captures reasoning stability: a model that produces diverse explanations for the same question may lack grounded reasoning.

\subsection{Metric 4: Hallucination Confidence Gap (HCG)}

HCG penalizes overconfident incorrect answers. Let $c = \exp\left(\frac{1}{n}\sum_{i=1}^{n} \log p_i\right)$ be the geometric mean of token probabilities (model confidence), and let $\text{F1}(a, g)$ be the token-level F1 score against the gold answer:
\begin{equation}
    \text{HCG} = \max(0, c - \text{F1}(a, g))
\end{equation}

HCG is zero when the model's confidence is justified by correctness, and positive when the model is overconfident relative to its accuracy.

\subsection{Composite FHI Score}

The Faithfulness-Hallucination Index combines all four metrics:
\begin{equation}
    \text{FHI} = \text{clip}\left(w_1 \cdot \text{AAS} + w_2 \cdot \text{CIS} + w_3 \cdot \text{ESS} - w_4 \cdot \text{HCG},\; 0,\; 1\right)
    \label{eq:fhi}
\end{equation}
where $w_1 = 0.25$, $w_2 = 0.35$, $w_3 = 0.25$, $w_4 = 0.15$ are empirically determined weights. CIS receives the highest weight as it provides the most direct causal evidence of faithfulness. HCG is subtracted as it represents a negative signal. A sample is classified as hallucinated if $\text{FHI} < \tau$, with threshold $\tau = 0.5$.
